% ----------------------------------------------------------------
% AMS-LaTeX definitions for a Honours Proposal *******************
% Template prepared by T.A. McWalter (May 2010)
% ----------------------------------------------------------------
\documentclass[a4paper,11pt,oneside]{article}

% PACKAGES -------------------------------------------------------

\usepackage{amsmath}
\usepackage{amsthm}
\usepackage{amsfonts}
\usepackage{amssymb}
\usepackage{ifpdf}
\usepackage{mathrsfs}
\usepackage[numbers]{natbib}
\usepackage{url}
\usepackage{xspace}

% GRAPHICS HANDLING ----------------------------------------------

\ifpdf
    \usepackage[pdftex]{graphicx}
    \DeclareGraphicsExtensions{.pdf, .jpg, .tif}
    \newcommand{\bb}{viewport}
\else
    \usepackage{graphicx}
    \DeclareGraphicsExtensions{.eps, .jpg}
    \newcommand{\bb}{bb}
\fi

% OTHER SET-UP ---------------------------------------------------

% Page dimensions

\setlength{\topmargin}{0.5cm}
\setlength{\headsep}{1cm}
\setlength{\textwidth}{14.0cm}
\setlength{\headheight}{14pt}
\setlength{\textheight}{22cm}
\setlength{\oddsidemargin}{1.3cm}
\setlength{\evensidemargin}{0cm}

% Line spacing

\newcommand{\setlinespacing}[1]
           {\renewcommand{\baselinestretch}{#1}\small\normalsize}


% MACRO DEFINITIONS ----------------------------------------------
\def\convK{\ensuremath{\kappa_t}\xspace}
\def\convKS{\ensuremath{\convK S_t}\xspace}
\def\convT{\ensuremath{\Omega_c}\xspace}

\def\putK{\ensuremath{P_t}\xspace}
\def\putT{\ensuremath{\Omega_p}\xspace}

\def\callK{\ensuremath{K_t}\xspace}
\def\callT{\ensuremath{\Omega_k}\xspace}

\def\defN{\ensuremath{R}\xspace}
\def\defS{\ensuremath{\eta}\xspace}

% PREAMBLE -------------------------------------------------------

\title{Honours Project Proposal:\\
Valuation of Convertible Bonds with Credit Risk}

\author{David Naylor\\[1cm]
Supervisor: Dr. Thomas McWalter\\[3cm]
Programme in Advanced Mathematics of Finance,\\
University of the Witwatersrand,\\
Johannesburg.\\[2cm]
\includegraphics[width=4cm]{../common/WitsColourLogo}\\[2cm]
}

% MAIN DOCUMENT --------------------------------------------------

\begin{document}

\maketitle
\thispagestyle{empty}

\setlinespacing{1.20}
\newpage

\setcounter{page}{1}

% ABSTRACT -------------------------------------------------------

\def\abstractname{\Large Abstract}
\begin{abstract}
Provide an abstract here.
\end{abstract}


% CONTENT --------------------------------------------------------

\section{Introduction}
This research project is proposed to price convertible bonds based on the paper by \citet{AFV03}.

Convertible bonds are a hybrid instrument available on financial markets.  The convertible bond is an instrument that is similar to a normal bond except the holder has the option to convert the bond into a specified number of shares.  The convertible bond typically also has embedded options whereby the issuer may buy-back the convertible bond for a specified price and whereby the investor can force the issue to repurchase the convertible bond.  In the event of a default by the issuer the bond could have partial recovery or total default.

More rigorously, as described by \cite{AKW08}, the convertible bond can be specified as having:
\begin{itemize}
 \item A nominal value of $N$.  Without loss of generality it is assumes the convertible bond always matures at par.
 \item A maturity at time $T$.
 \item A continuous conversion provision held by the investor to convert the bond into \convK shares (for a value of \convKS).  Conversion may only happen in time set \convT (typically $T \in \convT$).
 \item A continuous call provision held by the issuer to buy-back the bond for \callK.  The investor may opt to convert the bond instead of received the buy-back value.  The call may only happen in the time set \callT.
 \item A continuous put provision held by the investor to force early conversion of the bond for \putK.  The put may only happen in the time set \putT.
\end{itemize}

As assumed by \cite{AFV03}, the probably of default in time set $[t, dt]$ is $p\,dt$ where $p(S, t)$ is a deterministic hazard rate.  The following assumptions will be made about the behaviour of the bond in the event of default:
\begin{itemize}
 \item The bond has recovery ratio of \defN, for a recovery value of $\defN N$.
 \item The share price reduced in value by \defS so that the share price after default $S^+ = S^- (1 - \defS)$ where $S^-$ is the share price immediately before default.
 \item The investor has an option of either receiving:
 \begin{itemize}
  \item $\defN N$, or
  \item $\convK S (1 - \defS)$
 \end{itemize}
\end{itemize}

The stock price is assumed to follow a standard Wiener process of:
\begin{displaymath}
dS = \mu S dt + \sigma S dz
\end{displaymath}
with drift rate $\mu$, volatility of $\sigma$ and Wiener process increment of $dz$ \cite{AFV03}.

\section{Literature Review}

Provide a literature review here.

\section{Another Section if Necessary}

You can introduce other sections if necessary.

\section{Research Method and Aims}

End off by describing your proposed research methodology and aims.

% BIBLIOGRAPHY ---------------------------------------------------

\clearpage
\setlinespacing{1}
\addcontentsline{toc}{chapter}{Bibliography}
\bibliographystyle{../references/amsplnat}
\bibliography{../references/references.bib}
\end{document}
